\begin{thebibliography}{99}
\fontsize{9}{10.8}\selectfont
\setlength{\itemsep}{2.5pt}
\setlength{\parskip}{0pt}
\interlinepenalty=10000
\raggedright
\phantomsection
\addcontentsline{toc}{section}{Referencias}
\bibitem{feynman1948} R. P. Feynman, ``Space-Time Approach to Non-Relativistic Quantum Mechanics'', \emph{Reviews of Modern Physics} \textbf{20}, 367--387 (1948). \href{https://doi.org/10.1103/RevModPhys.20.367}{doi:10.1103/RevModPhys.20.367}.
\bibitem{mustovicary} B. Musto y J. Vicary, ``Quantum Latin Squares and Unitary Error Bases'', \emph{Quantum Information and Computation} \textbf{16}, 1318--1332 (2016). \href{https://doi.org/10.26421/QIC16.15-16-4}{doi:10.26421/QIC16.15-16-4}.
\bibitem{delascuevas2020} G. De las Cuevas, T. Drescher y T. Netzer, ``Quantum Magic Squares: Dilations and Their Limitations'', \emph{Journal of Mathematical Physics} \textbf{61}, 111704 (2020). \href{https://doi.org/10.1063/5.0022344}{doi:10.1063/5.0022344}.
\bibitem{delascuevas2023} G. De las Cuevas, T. Netzer e I. Valentiner-Branth, ``Magic Squares: Latin, Semiclassical, and Quantum'', \emph{Journal of Mathematical Physics} \textbf{64}, 022201 (2023). \href{https://doi.org/10.1063/5.0127393}{doi:10.1063/5.0127393}.
\bibitem{flm} I. B. Frenkel, J. Lepowsky y A. Meurman, \emph{Vertex Operator Algebras and the Monster}, Pure and Applied Mathematics, vol. 134, Academic Press, 1988.
\bibitem{borcherds} R. E. Borcherds, ``Monstrous Moonshine and Monstrous Lie Superalgebras'', \emph{Inventiones Mathematicae} \textbf{109}, 405--444 (1992). \href{https://doi.org/10.1007/BF01232032}{doi:10.1007/BF01232032}.
\bibitem{conwaysloane} J. H. Conway y N. J. A. Sloane, \emph{Sphere Packings, Lattices and Groups}, 3.\textsuperscript{a} ed., Grundlehren der mathematischen Wissenschaften 290, Springer, 1999. \href{https://doi.org/10.1007/978-1-4757-6568-7}{doi:10.1007/978-1-4757-6568-7}.
\bibitem{chenlamshimakura2016} H.-Y. Chen, C. H. Lam y H. Shimakura, ``\(\mathbb Z_3\)-orbifold construction of the Moonshine vertex operator algebra and some maximal \(3\)-local subgroups of the Monster'', prepublicación, arXiv:1606.05961 (2016), versión 2 (2017). \href{https://arxiv.org/abs/1606.05961}{arXiv:1606.05961}.
\bibitem{abelamyamada2017} T. Abe, C. H. Lam y H. Yamada, ``A remark on \(\mathbb Z_p\)-orbifold constructions of the Moonshine vertex operator algebra'', prepublicación, arXiv:1705.09022 (2017), versión 4. \href{https://arxiv.org/abs/1705.09022v4}{arXiv:1705.09022}.
\bibitem{kmconwaynorton1979} J. H. Conway y S. P. Norton, ``Monstrous Moonshine'', \emph{Bulletin of the London Mathematical Society} \textbf{11}, 308--339 (1979). \href{https://doi.org/10.1112/blms/11.3.308}{doi:10.1112/blms/11.3.308}.
\bibitem{Cantor1878} G. Cantor, ``Ein Beitrag zur Mannigfaltigkeitslehre'',
\emph{Journal für die reine und angewandte Mathematik} \textbf{84},
242--258 (1878). \href{https://doi.org/10.1515/crelle-1878-18788413}{doi:10.1515/crelle-1878-18788413}.
\bibitem{Hilbert1902} D. Hilbert, ``Mathematical Problems'',
\emph{Bulletin of the American Mathematical Society} \textbf{8},
437--479 (1902). \href{https://doi.org/10.1090/S0002-9904-1902-00923-3}{doi:10.1090/S0002-9904-1902-00923-3}.
\bibitem{Godel1938} K. Gödel, ``The Consistency of the Axiom of Choice and
of the Generalized Continuum-Hypothesis'', \emph{Proceedings of the National
Academy of Sciences} \textbf{24}, 556--557 (1938).
\href{https://doi.org/10.1073/pnas.24.12.556}{doi:10.1073/pnas.24.12.556}.
\bibitem{Cohen1963} P. J. Cohen, ``The Independence of the Continuum
Hypothesis'', \emph{Proceedings of the National Academy of Sciences}
\textbf{50}, 1143--1148 (1963).
\href{https://doi.org/10.1073/pnas.50.6.1143}{doi:10.1073/pnas.50.6.1143}.
\bibitem{Cohen1964} P. J. Cohen, ``The Independence of the Continuum
Hypothesis, II'', \emph{Proceedings of the National Academy of Sciences}
\textbf{51}, 105--110 (1964).
\href{https://doi.org/10.1073/pnas.51.1.105}{doi:10.1073/pnas.51.1.105}.
\enlargethispage{3\baselineskip}
\bibitem{Jech2003} T. Jech, \emph{Set Theory}, 3.\textsuperscript{a} ed.,
Springer Monographs in Mathematics, Springer, 2003; capítulo 13,
``Constructible Sets'', pp. 175--200.
\href{https://doi.org/10.1007/3-540-44761-X}{doi:10.1007/3-540-44761-X}.
\end{thebibliography}
